\subsection{A complete nonzero-first-coordinate global section calculation}

\begin{statement}{E.49}{Proved}{Every global section of the inherited quotient with first Kummer class three}
Use the actual arithmetic quotient $H$ and origin
of E.38 and E.46. Choose positive compatible real
roots of $3$, put $\kappa=\kappa_3(3)$, and fix
the first cocycle $A=3\kappa$. Define the explicit
set of rational primes
\[
 \mathscr P=\{p\ne3:\ 3\nmid
          \operatorname{ord}_{\mathbb F_p^\times}(27\bmod p)\}.
 \tag{E49.1}
\]
For $p=2$ the order is one. Equivalently the
three-primary part of the order of $3\bmod p$
has order at most three. In particular every
prime $p\equiv2\pmod3$, including two, belongs
to $\mathscr P$.

Let
\[
 \mathcal R(\mathscr P)=
 \{(b_p)_p\in\prod_{p\in\mathscr P}\mathbb Z_3:
        \text{for each }r,\ b_p\in3^r\mathbb Z_3
                     \text{ for all but finitely many }p\}.
 \tag{E49.2}
\]
It carries its multiplication-by-$3^r$
inverse-limit topology. The global section
classes of $H$ with first class $[3\kappa]$
have complete coordinates
\[
 \boxed{\quad
   (b_3,(b_p)_{p\in\mathscr P},z,f)
   \in \mathbb Z_3\times\mathcal R(\mathscr P)
                  \times\mathbb Z_3\times\mathbb Z_3[[Z]] .
 \quad}
 \tag{E49.3}
\]
The proof constructs a homeomorphism after
fixing one specified norm lift at each prime
of $\mathscr P$. The toric map scaled as in
(E46.7) is exactly
\[
 (b_3,(b_p),z,f)\longmapsto
 \left(\kappa,\,
         b_3\kappa+\sum_{p\in\mathscr P}b_p\kappa_3(p)\right).
 \tag{E49.4}
\]
Thus its image consists exactly of completed
rational Kummer classes whose valuations at
primes outside $\{3\}\cup\mathscr P$ vanish.
Every fiber has the explicit compact
coordinates $\mathbb Z_3\times\mathbb Z_3[[Z]]$.
The entire section group, with the origin
$g\mapsto(A_g,0)$, is torsion-free.

Here are formulas for the compact parameters.
Write
\[
 \Lambda=\mathbb Z_3[[U]],\quad T=1+U,\quad
 W=(T-T^{-1})/2,\quad
 \eta=1-(T^{-1}+T^{-2})/2.
 \tag{E49.5}
\]
The proof selects actual conjugate Frobenii
$\tau,\tau'$ over seven, still a free basis
of $P=G_{\mathbb Q(\zeta_3),\{3\}}^{(3)}$,
with $\kappa(\tau)=1$ and
$\kappa(\tau')=-1$. Every section class
unramified outside three has a unique
second-cocycle representative with
\[
 \begin{split}
 v(c)&=0,\\
 v(\tau)&=3b_3+z\eta+T^3W f(W^2),\\
 v(\tau')&=-v(\tau)(T^{-1}).
 \end{split}
 \tag{E49.6}
\]
These values define the cocycle on the
actual free group and its extension by
complex conjugation. All coefficients
of the power series $f$ are arbitrary
elements of $\mathbb Z_3$.

There is also a complete unit and valuation
interpretation of (E49.3). For $n\ge1$ put
\[
 \begin{gathered}
 \alpha_n=3^{1/3^{n-1}},\qquad
 \beta_n=\alpha_n\zeta_{3^n},\\
 F_n=\mathbb Q(\alpha_n),\qquad
 L_n=\mathbb Q(\beta_n),\qquad
 E_{A,n}=F_n\times L_n .
 \end{gathered}
 \tag{E49.7}
\]
The fields have degrees $3^{n-1}$ and
$2\cdot3^{n-1}$, and their tower maps
are $\alpha_{n+1}^3=\alpha_n$ and
$\beta_{n+1}^3=\beta_n$. The inverse norm
limit of their completed ordinary global
units is explicitly
\[
 \varprojlim_N
   \left(\widehat{\mathcal O_{F_n}^{\times}}_3
                \times\widehat{\mathcal O_{L_n}^{\times}}_3\right)
 \ \xrightarrow{\sim}\ \mathbb Z_3[[Z]].
 \tag{E49.8}
\]
The inverse map $\epsilon(f)$ is the exact
Kummer--Shapiro norm system of the
unramified cocycle with $\tau$-value
$T^3W f(W^2)$.

Define two norm-coherent three-unit systems
$\xi_F,\xi_L$ by the unramified cocycles
with $\tau$-values $1$ and $1-\eta$,
respectively, in $\Lambda_A$.
The combination (E49.9) will satisfy
the additional $J_A$ condition.
They have respective
three-prime valuation pairs $(1,0)$
and $(0,1)$ in $F_n,L_n$, and both
have level-zero norm $3$.
For each $p\in\mathscr P$, fix a norm
system $\pi_p$ whose divisor at each
level is the unique compatible degree-one
prime above $p$, with zero valuations at
every other prime, including three.
The proof constructs all these systems
using surjective unit norm maps. Their
level-zero norms are exactly $p$.
The complete norm-to-section map in
(E49.3) is then
\[
       \epsilon(f)\,
       \xi_F^{\,3b_3+z}\xi_L^{-z}
              \prod_{p\in\mathscr P}\pi_p^{\,3b_p}.
 \tag{E49.9}
\]
Every product and power is taken in
the completed multiplicative groups
of E.46. No hypothesis on an unevaluated
class group or unit norm map remains.

The geometric map is precisely $N\to H$
with kernel $[I_y,I_y]$; its factorizations
through the intrinsic metabelian and marked
groups have the kernels (E47.8). Statement
(E49.3) evaluates the entire indicated
global section fiber of this target.

\emph{Manuscript: \S51.16, folios733v--736v,
scans172--178, (371)--(379), Part E(c),(d).
The calculation retains the actual arithmetic
origin, first-coordinate cocycle, and
completed-group-ring topology.}
\end{statement}
\begin{proof}
\emph{The actual arithmetic basis and its normalization.}
Let $K=\mathbb Q(\zeta_3)$ and retain the
positive real roots defining $\kappa$.
Then $\kappa(c)=0$. Start with the
actual Frobenii $\tau_0,\tau'_0=c\tau_0c$
of E.37, taking the prime over seven
at which $\zeta_3$ reduces to two.
The cubic residue calculation
$3^{(7-1)/3}=2=\zeta_3$ gives
$t=\kappa(\tau_0)\equiv1\pmod3$.
The cocycle identity gives
$\kappa(\tau'_0)=-t$.

Put $\rho_0=\tau'_0\tau_0^{-1}$.
Its cyclotomic character is one and
$A(\rho_0)=-6t$. The three-adic power
\[
        h=\rho_0^k,\qquad
             k=(t^{-1}-1)/12\in\mathbb Z_3
 \tag{E49.10}
\]
is therefore defined, and
$A(h)=-6tk$. Set
$\tau=h\tau_0h^{-1}$ and $\tau'=c\tau c$.
The exact cocycle formula for conjugation
gives
\[
 A(\tau)=3t+(1-7)(-6tk)=3,\qquad
                   A(\tau')=-3.
 \tag{E49.11}
\]
Both elements remain actual Frobenii at
conjugate prolongations over seven.
Their Frattini classes are the unchanged
classes of $\tau_0,\tau'_0$; the Frattini
criterion in the free rank-two group
therefore proves that they remain a
free basis. This supplies the arithmetic
map behind the normalization.

Use the Nielsen basis
$\sigma=\tau$, $\rho=\tau'\tau^{-1}$.
It satisfies
\[
 c\sigma c=\rho\sigma,\qquad c\rho c=\rho^{-1}.
 \tag{E49.12}
\]
On the twisted coefficient module
$\Lambda_A$ of E.46, the actions are
the explicit operators
\[
 \sigma u=7T^3u(T^7),\qquad
 \rho u=T^{-6}u,\qquad
 Cu=c u=-u(T^{-1}).
 \tag{E49.13}
\]
Their compositions verify (E49.12)
on the coefficient module as well.

\emph{Finite coefficients and the full arithmetic group.}
We record the exact comparison used below.
For a finite three-primary module $M$
whose action over $K$ factors through $P$,
inflation induces an isomorphism
\[
 H^i(P,M)\ \longrightarrow\
                  H^i(G_{K,\{3\}},M)
             \qquad(i\ge0).
 \tag{E49.14}
\]
For degree one, a cocycle gives a
homomorphism into $M$ semidirect its
finite three-group of arithmetic actions.
This target is a three-group, so the
homomorphism factors through the actual
maximal pro-three quotient $P$.
Coboundaries use the same elements of
$M$ on both sides, proving injectivity
as well as surjectivity. Degree zero
is immediate.

For higher degrees, every simple factor
of $M$ for its finite three-group action
is the trivial module $\mathbb F_3$.
The established calculation of E.34
gives $H^2(G_{K,\{3\}},\mathbb F_3)=0$;
the arithmetic cohomological dimension
bound there gives zero in higher degrees.
Induction along a finite composition
series gives the same vanishing for
$M$. The free group $P$ has
three-cohomological dimension one.
This proves the asserted comparison
in every degree for precisely these
coefficient modules.

In particular take the finite images
$\Lambda_r,J_r$ in
$(\mathbb Z/3^r)[\mathbb Z/3^r]$.
Their action over $K$ has three-power
image: the positions transform by
$i\mapsto A_g+\chi_gi$ with
$A_g\in3\mathbb Z_3$ and
$\chi_g\in1+3\mathbb Z_3$, and the
scalar twist has three-power image
as well. Thus (E49.14) applies.
Restriction to $K$ followed by
corestriction multiplies by two, so
the corresponding degree-two and
degree-three groups over
$G_{\mathbb Q,\{3\}}$ also vanish.
Its degree-one groups with these
finite coefficients are finite:
over $P$ a cocycle is specified by
two values in a finite module, and
descent adds only the finite group
of order two.

Taking the cofinal system $r\to\infty$
therefore retains vanishing in degree
two. In particular the inverse systems
of unramified degree-one groups are
Mittag--Leffler. The continuous-cochain
comparison is the finite-coefficient
comparison proved in E.44 and
\cite[Proposition 1.6 and Theorem 2.2]{E_JannsenContinuous}.
No comparison with a pro-three quotient
for an unrelated coefficient action
is being assumed.

\emph{An explicit normal form for all unramified classes.}
Normalize a second cocycle at $c$ to
zero using division by two. It is then
freely specified by
$u=v(\tau)\in\Lambda$ and
$v(\tau')=Cu$. The remaining
coboundaries are
\[
     Dm=7T^3m(T^7)-m,\qquad
     m\in\Lambda^C
        =\{m:m(T^{-1})=-m(T)\}.
 \tag{E49.15}
\]
The first-degree comparison in (E49.14)
and the free property of $P$ prove
that these are exactly the global
classes unramified outside three.
For $J_A$ one takes $u\in J$;
the residual fixed lattice is the
same, because an element of
$\Lambda^C$ has augmentation zero.

Here is a reduction algorithm entirely
inside the integral completed ring.
Put $n=T^{-6}-1$. The exact identity
\[
       n=-T^{-6}(T^3+1)(T^3-1)
 \tag{E49.16}
\]
identifies its ideal with $(T^3-1)$:
the extra factor is a unit. Consequently
$\Lambda/(n)=\mathbb Z_3[C_3]$, with
the specified representatives $1,T,T^2$.
For $u\in\Lambda$, let its image be
$u_0+u_1T+u_2T^2$ and set
\[
 \begin{gathered}
       a=\operatorname{aug}u,\qquad z=-(u_1+u_2),\\
       u_* =u-a-z\eta,\qquad
       m=\frac{Cu_*-T^{-6}u_*}{T^{-6}-1}.
 \end{gathered}
 \tag{E49.17}
\]
The numerator has zero image in
$\mathbb Z_3[C_3]$: the image of
$u_*$ is antisymmetric under inversion.
Thus the division is integral.
Multiplication by $n$ is injective
in the domain $\Lambda$.
Applying $C$ to the numerator shows
$m=Cm$; explicitly
$C(nm)=-T^6n\,Cm$ and the numerator
equals $-T^{-6}$ times its image under
$C$. Therefore (E49.15) permits
subtracting $Dm$.

The resulting
$v=u_*-Dm$ satisfies
$Cv=T^{-6}v$. Thus $T^{-3}v$
is antisymmetric under $T\mapsto T^{-1}$.
The variable $W$ in (E49.5) has
linear coefficient one in $U$.
Its inverse substitution is
$T=W+\sqrt{1+W^2}$, with the
unique square root of constant
term one. Hence
$\Lambda=\mathbb Z_3[[W]]$, inversion
sends $W$ to $-W$, and
\[
              v=T^3W f(W^2)
                    \quad\text{for a unique }f\in\mathbb Z_3[[Z]].
 \tag{E49.18}
\]
This gives the normal form
$a+z\eta+T^3W f(W^2)$.

It is unique modulo (E49.15).
Augmentation of $Dm$ is zero, so
it cannot change $a$. The
$\rho$-value of a cocycle with
$\tau$-value $u$ is
$Cu-T^{-6}u$. Its image modulo
$n$ detects $z$ and is unchanged
by the coboundary, whose
$\rho$-value is $nm$. Once $a,z$
are zero, a representative of
the remaining form has
$\rho$-value zero; a coboundary
with this property has $nm=0$,
hence $m=0$. The power series
is consequently unchanged.
All the division maps just used
are continuous on their closed
principal ideals, since they
invert continuous bijections
from a compact ring. This also
proves the normal-form topology.
For $u\in J$ its augmentation
is $a=3b_3$, giving (E49.6).

\emph{The radical fields and surjectivity of their unit norms.}
The rational polynomial factors as
\[
 X^{3^n}-27
  =(X^{3^{n-1}}-3)
       (X^{2\cdot3^{n-1}}+3X^{3^{n-1}}+9).
 \tag{E49.19}
\]
The first factor is Eisenstein at
three. For the second,
$\beta_n^{3^{n-1}}=3\zeta_3$,
so $\beta_n$ contains $K$.
Normalize the valuation of $K$
at $\lambda=1-\zeta_3$ by
$v_\lambda(3)=2$.
The valuation $2/3^{n-1}$ of
$\beta_n$ forces ramification
degree at least $3^{n-1}$ over
$K$, since its numerator is
prime to three. This proves
the full degree, irreducibility
of the second factor, and the
unique totally ramified prime
above three. The same argument
applies to $M_n=K F_n$.

Both $L_{n+1}/L_n$ and
$M_{n+1}/M_n$ are cyclic cubic:
the base contains $\mu_3$ and
the next generator has cube
equal to the preceding one.
They are unramified away from
three, as seen from the
polynomial discriminants,
and have exactly one totally
ramified prime. Starting from
$L_1=M_1=K$, the ambiguous
class-number argument of E.47
therefore gives
$3\nmid h(L_n)h(M_n)$ at every
level. Ideal extension followed
by norm from $M_n$ multiplies
$\operatorname{Cl}(F_n)$ by
two, so $3\nmid h(F_n)$ also.

The strongly ambiguous formula
in \cite[Theorem 1, formula (2)]{E_LemmermeyerAmbiguous}
has the especially useful
consequence, for each of these
cyclic cubic layers $L'/L$,
\[
       |\operatorname{Am}_{\rm st}(L'/L)|
         =\frac{h(L)}
                   {[\mathcal O_L^\times:
                         N\mathcal O_{L'}^\times]}.
 \tag{E49.20}
\]
The denominator is a power of
three, because the cube of
each base unit is a norm.
It divides the prime-to-three
integer $h(L)$, and is
therefore one. Actual ordinary
unit norms in the $L_n$ and
$M_n$ towers are onto.
For the $F_n$ tower, project
the latter norm map to the
fixed submodule of complex
conjugation by $(1+c)/2$.
The fixed submodule of the
completed units of $M_n$
is exactly the completed
units of $F_n$: norm followed
by inclusion is $1+c$,
and two is invertible.
This proves surjectivity of
the completed ordinary unit
norms for both factors of
(E49.7).

\emph{The ordinary-unit norm limit in exact cochain coordinates.}
Apply finite Shapiro over
$G_{\mathbb Q,\{3\}}$ to
the positional modules.
Every field in (E49.7)
is unramified outside three.
Kummer theory identifies
$H^1(G_{F,\{3\}},\mathbb Z_3(1))$
with the completed three-units:
the inverse Tate module of
the finite $S$-class group
is zero. In this tower its
three-primary class group
is already zero by the
preceding paragraph. The
norm comparison of E.46
therefore identifies the
unramified cohomology
normal form with the entire
inverse norm limit of the
completed three-units.

To obtain the ordinary units,
compute the two three-prime
valuations at level one.
There $E_{A,1}=\mathbb Q\times K$.
Write $u\bmod(T^3-1)=u_0+u_1T+u_2T^2$.
Shapiro at the rational root
gives valuation $u_0$,
because $\kappa(\tau)=1$.
At the unit-root orbit,
the two $K$-cocycle values
are $u_1$ and $-u_2$.
The torsion Kummer class
of $\zeta_3$ has values
$2,2$, since it is represented
over $K$ by $(\chi-1)/3$.
The relation
$\lambda^2=-3\zeta_3$
gives
$2[\kappa_3(\lambda)]
  =[\kappa]+[\kappa_3(\zeta_3)]$.
Thus the two values of its
class have difference one.
A coboundary changes both
values by the same $6m$.
The valuation at $\lambda$
of the unit-orbit class is
therefore exactly $u_1+u_2$.

For the normal form just
proved, these two valuations
are
\[
                 (a+z,-z).
 \tag{E49.21}
\]
Indeed $T^3W f(W^2)$ reduces
to an antisymmetric vector
modulo $T^3-1$, while
$\eta$ reduces to
$1-(T+T^2)/2$.
The residue degrees at the
unique three-prime in each
tower layer are one, so norm
compatibility makes these
valuations constant in the
tower. Consequently the
ordinary-unit norm limit
is exactly $a=z=0$.
This proves (E49.8), including
its explicit map and compact
topology. It also proves the
asserted valuations and total
norms of $\xi_F,\xi_L$.

\emph{The surviving primes and the exact residue lattice.}
At $p\ne3$ the coefficient
action is unramified.
For a chosen actual
arithmetic Frobenius $F_p$
it acts, after removing the
Tate twist, by pushforward
under
\[
             \phi_p(x)=A_{F_p}+px
                         \quad(x\in\mathbb Z_3).
 \tag{E49.22}
\]
Its only possible fixed
point is
$x_p=-A_{F_p}/(p-1)$.
An integral fixed point
is exactly a compatible
system of $3^n$th roots
of $27$ in $\mathbb Q_p$,
using the root labels
of E.46. Cubing is an
automorphism on the
principal units at
$p\ne3$. Thus those
compatible roots exist
exactly when the residue
of $27$ has order prime
to three, proving (E49.1).
For $p\equiv2\pmod3$,
$p-1$ is a three-adic
unit and the fixed point
is always integral.

The invariant integral
measure lattice for
$\phi_p$ is exactly
$\mathbb Z_3\delta_{x_p}$
when this point is integral,
and is zero otherwise.
Here is the integral
argument. With an integral
center, translate it to
zero and use multiplication
by $p^2$. On every nonzero
annulus its clopen orbits
are free $\mathbb Z_3$
orbits, as in E.45.
An invariant integral
measure on a free orbit
has its mass divisible
by arbitrarily large
powers of three, hence
vanishes; equivalently
$S-1$ is injective on
$\mathbb Z_3[[S-1]]$.
Only the atom at the
fixed point remains.
When the center is not
integral, necessarily
$p\equiv1\pmod3$.
Put $s=v_3(p-1)$ and
$t=v_3(A_{F_p})<s$.
The identity
$v_3(\phi_p^n(x)-x)=t+v_3(n)$
shows that the clopen
sets $i+3^t\mathbb Z_3$
are free orbits. The
same measure argument
gives zero. Intersecting
the invariant lattice
with $J$, whose masses
are divisible by three,
gives precisely
\[
 (J(-1))^{F_p}=
 \begin{cases}
      3\mathbb Z_3\delta_{x_p},&p\in\mathscr P,\\
      0,&p\notin\mathscr P,\ p\ne3.
 \end{cases}
 \tag{E49.23}
\]

This also computes the
valuation inverse system
on the norm side. At
each finite positional
level, the Frobenius
orbits correspond to
the primes of $E_{A,n}$
above $p$. The invariant
vector $\sum_{i\in O}\delta_i$
represents valuation one
at that prime. Positional
pushforward multiplies it
by the relative residue
degree, exactly as does
the ideal norm.
The fixed point $x_p$
therefore gives the unique
compatible degree-one
prime chain. Every other
valuation is forced to
zero by its unbounded
three-power residue degrees.
At three there are the
two totally ramified
chains of residue degree
one in (E49.7).
The complete valuation
inverse limit is consequently
$\mathbb Z_3^2\times\mathcal R(\mathscr P)$.
Restrictedness is exact:
the mass at level zero
forces the coefficients
at distinct $p$ to tend
to zero modulo every
power of three, and that
condition is sufficient
in every positional level.

\emph{Residue surjectivity over the actual arithmetic group.}
For clarity the same
surjectivity can be proved
before choosing elements
of the radical fields.
Localization, purity, and
the group-cohomology
comparison used in E.44
give, for each $J_r$, the
exact sequence
\[
 \begin{gathered}
 0\longrightarrow
 H^1(G_{\mathbb Q,\{3\}},J_r)
 \longrightarrow H^1(G_{\mathbb Q},J_r)\\
 \longrightarrow
   \bigoplus_{p\ne3}(J_r(-1))^{F_p}
 \longrightarrow0 .
 \end{gathered}
 \tag{E49.24}
\]
The last zero is the
finite-coefficient
degree-two vanishing
proved above. The
primary maps are
\cite[II.0.1(f), II.1.7(b), II.2.9]{E_MilneADT}.
The left groups are
finite, so their first
inverse-limit derivative
vanishes. The degree-zero
finite groups needed to
compare continuous
degree-one cohomology
also satisfy
Mittag--Leffler.
Taking the inverse limit
therefore preserves
surjectivity and gives
exactly the restricted
fixed measures in
(E49.23), without
taking a closure or
discarding an inverse-limit
term. The tame generator
can be normalized by
$\kappa_3(p)(\tau_p)=1$;
the residue $3b_p\delta_{x_p}$
then has scaled toric
valuation $b_p$.

\emph{Constructing all norm systems and all section coordinates.}
There is an equally
explicit construction
using the fields.
At each finite level,
the completed divisor
map is onto: the class
groups of $F_n,L_n$
have order prime to
three. If an ideal
has class of order $h$
prime to three, choose
a generator of its
$h$th power and take
its $h^{-1}$ power
in the three-completion.
This produces its
valuation vector exactly.
The kernel of the
completed divisor map
is the completed
ordinary unit group.
For example this follows
by splitting the
sequence onto the
free principal-divisor
image, then using its
prime-to-three index
in the full divisor
group.

For $p\in\mathscr P$,
choose a completed
element with divisor
the prime at each
level of its degree-one
chain. If its norm
differs from the
element already chosen
downstairs, the ratio
is an ordinary unit.
The surjective unit
norms proved in
(E49.20) correct this
ratio. Induction
constructs $\pi_p$
with exactly the
specified divisors.
At level zero its
norm is $p$, because
the completion of
$\mathbb Q$'s ordinary
unit group is zero.
This construction
uses actual norm
equations, not an
assumed saturation
of the divisor image.

Any compatible norm
system has unique
valuations $(d_F,d_L)$
at three and $(c_p)_p$
on these other chains.
After dividing by
$\xi_F^{d_F}\xi_L^{d_L}
\prod\pi_p^{c_p}$,
it is an ordinary-unit
system, and has the
unique form $\epsilon(f)$
by (E49.8).
The infinite product
converges because
$c_p\to0$ three-adically;
at every coefficient
exponent only finitely
many factors contribute.
Its level-zero norm
has valuations
$d_F+d_L$ at three
and $c_p$ at $p$.
The cube condition
of E.46 is exactly
\[
       d_F+d_L\in3\mathbb Z_3,\qquad
              c_p\in3\mathbb Z_3\quad(p\in\mathscr P).
 \tag{E49.25}
\]
Set $d_F=3b_3+z$,
$d_L=-z$, and $c_p=3b_p$.
This proves (E49.9),
its uniqueness, and
the complete coordinates
(E49.3)--(E49.4).

All valuation maps are
continuous in the
power-class topologies.
The products just
constructed are
continuous for those
topologies, and the
remaining unit map
and its inverse are
continuous by the
compact normal-form
calculation. Thus the
classification is a
homeomorphism. Changes
of the chosen systems
$\pi_p$ multiply them
by explicit ordinary-unit
norm systems, and alter
the $f$ coordinate by
the corresponding
convergent linear sum.
They do not change
the toric criterion
or any fiber.
Every parameter group
in (E49.3) is
three-torsion-free,
proving the stated
torsion assertion.
\end{proof}
